\subsubsection{Current cones and endpoint equations}
\label{sec:warped-current-endpoint-equations}

Retain $X=B\times C$, $\theta=b\kappa+(ca+e)w$ with real parameters $b,c,e$, and the probability $\mu=\sigma\wedge w$ of \cref{sec:warped-gauduchon-completion}. Localized closed tests determine its current cone as well as its smooth cone. At a finite unit-volume end, the reciprocal coefficient is integrable with an exponent greater than one, so the base equation still has a continuous solution. Write $h=\kappa+w$.

\Needspace{14\baselineskip}
\begin{proposition}\label{prop:warped-current-cones}\index{Cone duality!warped product currents}
In this real span, the pseudoeffective and nef cones both equal
\[
 \{b\geq0,\ e\geq0,\ e+cM\geq0\}.
\]
A background admits a strict positive current representative if and only if $b>0$, $e>0$ and $e+cM>0$. Thus its strict-current cone equals the smooth-positive cone $\mathscr C$ in \cref{thm:warped-gauduchon-volume}.
\end{proposition}
\begin{proof}
Put $L=ca+e$. Suppose an arbitrary quasi-plurisubharmonic $u$ satisfies $\theta+\ddc u\geq\delta h$, where $\delta\geq0$. The closed semipositive test $\kappa^{m-1}\wedge w$ and distributional Stokes give $b\geq\delta$. For every smooth nonnegative top-degree form $\eta$ on $B$, its pullback to $X$ is closed by dimension, and
\[
 \int_X(\theta+\ddc u-\delta h)\wedge\eta
   =\int_B(L-\delta)\eta\geq0.
\]
If $L<\delta$ at a point, a nonzero test supported in that open set contradicts this inequality. Hence $L\geq\delta$ everywhere. For $\delta=0$ these are precisely $b\geq0$, $e\geq0$, $e+cM\geq0$; conversely they make $\theta$ semipositive. Such forms are nef because $\theta+\varepsilon h>0$ for every $\varepsilon>0$. Conversely, applying the same tests to smooth positive representatives of $\theta+\varepsilon h$ and sending $\varepsilon$ to zero gives those inequalities. For $\delta>0$ they force the three strict inequalities, and the converse follows because $\theta$ itself is then strictly positive. No restriction of $u$ to an extreme fiber is needed.
\end{proof}

\Needspace{14\baselineskip}
\begin{proposition}\label{prop:warped-continuous-endpoint-equation}\index{Monge--Ampere equation@Monge--Amp\`ere equation!finite warped boundary}\index{Regularity!continuous solution without C11@continuous solution without $C^{1,1}$}
Under the extreme-set hypotheses of \cref{lem:warped-endpoint-length}, let $\theta=b\kappa+L_0w$ be either finite face $\theta_+(b)$ or $\theta_-(b)$ of \cref{thm:warped-unit-volume-completion}, with $b>0$. There is a continuous $\theta$-plurisubharmonic potential $u$, pulled back from $B$, such that
\begin{equation}\label{eq:warped-continuous-endpoint-equation}
 (\theta+\ddc u)^n=\mu.
\end{equation}
The product is the ordinary bounded-potential product. No solution of this equation is globally $C^{1,1}$.
\end{proposition}
\begin{proof}
Let $f=a$ at the first face and $f=1-a/M$ at the second, and let $r>1$ be the corresponding codimension-to-order ratio. The coefficient $L_0$ is a positive constant multiple of $f$. In each extreme tube the integral of $f^{-p}$ has the same convergence condition as
\[
 \int_0^{s_0}s^{\operatorname{codim}_{\mathbb R}Z_i-1-k_ip}\,\mathrm{d}s,
\]
where $s$ is normal radius and $s_0>0$ is small. Thus $1/L_0\in L^p(B,\sigma)$ for any fixed $1<p<r$. The zero set has smooth volume zero, and the density can be assigned any finite value there. The face formulas give
\[
 H_0:=\int_B\frac{\sigma}{L_0}=nb^m.
\]
Apply \cref{thm:kn-continuous} on the positive closed background $b\kappa$ to obtain a continuous $\phi$ with $A=b\kappa+\ddc\phi\geq0$ and $A^m=\lambda\sigma/L_0$ for a scalar $\lambda>0$. The bounded closed mass identity in \cref{thm:mass-invariance-characterization} gives $\int_BA^m=b^m$, hence $\lambda=b^m/H_0$. Set $u=\phi\circ p_B$, where $p_B:X\to B$ is projection. Its solution current is $A+L_0w\geq0$.

The pullback identity of \cref{lem:bounded-products-submersion} gives $A^{m+1}=0$ on $X$, and $(L_0w)^2=0$. The multilinear expansion in \cref{prop:bounded-products} therefore gives
\[
 (A+L_0w)^n=nL_0 A^m\wedge w
       =\frac{nb^m}{H_0}\sigma\wedge w=\mu.
\]
The cancellation $L_0(\sigma/L_0)=\sigma$ holds as measures, since the zero set has smooth volume zero and $A^m$ has the displayed absolutely continuous density.

Suppose a solution $q$ were globally $C^{1,1}$. Since $\mu/h^n$ is smooth and strictly positive on compact $X$, \cref{prop:c11-ellipticity} would give $\theta+\ddc q\geq\delta h$ for some $\delta>0$. This contradicts \cref{prop:warped-current-cones}, because $L_0$ vanishes. The argument applies to every solution, not only to the pulled-back one.
\end{proof}

\subsubsection{Limits of the base potentials}
\label{sec:warped-endpoint-potential-limits}

The normalized potentials used to construct the interior balanced metrics have a unique continuous limit at each finite face. The compactness and uniqueness needed for this conclusion concern the closed base equation.

\Needspace{15\baselineskip}
\begin{proposition}\label{prop:warped-endpoint-potential-limits}\index{Gauduchon metric!boundary potential limits}\index{Compactness!warped endpoint potentials}
Under the hypotheses of \cref{thm:warped-unit-volume-completion}, let $\theta_j=b_j\kappa+L_jw\in\mathscr S$ converge in coefficients to a finite face $\theta_0=b\kappa+L_0w$. Normalize the smooth base potentials in \cref{thm:warped-gauduchon-volume} by $\sup_B\phi_j=0$ and put $u_j=\phi_j\circ p_B$. Then $u_j$ converge uniformly on $X$ to a continuous pulled-back solution $u_0$ of \eqref{eq:warped-continuous-endpoint-equation}. For every $1\leq k\leq n$,
\[
 (\theta_j+\ddc u_j)^k\rightharpoonup
             (\theta_0+\ddc u_0)^k.
\]
It is the unique normalized bounded solution pulled back from the base.
\end{proposition}
\begin{proof}
Write $v_j=c_j/e_j$ and $\psi_j=\phi_j/b_j$. The normalized closed base equation is
\[
 (\kappa+\ddc\psi_j)^m=F_{v_j}\sigma,
 \qquad F_v=\frac1{Q(v)(1+va)},\qquad\sup_B\psi_j=0.
\]
At the first face, $v_j\to\infty$. Set $\epsilon=1/v$ and $H_\epsilon=\int_B\sigma/(a+\epsilon)$. Then $F_v=[H_\epsilon(a+\epsilon)]^{-1}$ and $H_\epsilon\to H_a$. At the second face, $v_j\to-1/M$. With $f=1-a/M$ and $t=v+1/M>0$, the denominator is $Q(v)(f+ta)$ and $Q(v)\to H_-$. Thus the respective limits are
\[
 F_+=\frac1{H_a a},\qquad F_-=\frac1{H_-(1-a/M)}.
\]
Choose $1<p<r$, with $r>1$ the ratio for that face. Tubular integration as in \cref{prop:warped-continuous-endpoint-equation} gives $a^{-1}\in L^p(B,\sigma)$ or $f^{-1}\in L^p(B,\sigma)$, respectively. Near the first end, $F_v\leq C/a$; near the second, $F_v\leq C/f$. The extreme sets have smooth volume zero. Dominated convergence therefore gives $F_{v_j}\to F_\pm$ in $L^p(B,\sigma)$, and each limiting density has integral one. Since $\sigma/\kappa^m$ is smooth and strictly positive on compact $B$, the equation densities also have a common $L^p(B,\kappa^m)$ bound.

Apply \cref{thm:lp-uniform-compactness} to the fixed background $\kappa$. Every subsequence of $\psi_j$ has a uniformly convergent further subsequence. Bounded-product continuity identifies its limit with a normalized solution of $(\kappa+\ddc\psi_0)^m=F_\pm\sigma$. Closedness of $\kappa$ gives the Guan--Li condition, so \cref{uni-thm:GL-intro} makes this solution unique. Every cluster limit is consequently the same, which proves whole-sequence uniform convergence. Since $b_j\to b>0$, we obtain $u_j\to u_0=(b\psi_0)\circ p_B$ uniformly.

On each coordinate ball, coefficient convergence of the backgrounds gives a common smooth quadratic correction making all $u_j+q$ plurisubharmonic. These corrected potentials are uniformly bounded and converge uniformly. The capacity-continuity statement in \cref{sec:psh-local-capacity} gives weak convergence of all their bounded products, with uniform local mass bounds. Expanding off the correction as in \eqref{eq:bounded-expansion}, and using smooth convergence of the remaining background factors, gives the displayed current convergence. In top degree the interior products all equal $\mu$, so the limit solves the endpoint equation.

For any bounded pulled-back solution $v=\chi\circ p_B$ with $\sup_B\chi=0$, the horizontal current $A_\chi=b\kappa+\ddc\chi$ on $B$ is positive: pair its pulled-back current with positive base tests and $w$. The product expansion gives $nL_0A_\chi^m=\sigma$. Off the extreme set this determines $A_\chi^m=\sigma/(nL_0)$. That density already has total mass $H_0/n=b^m$, while closed bounded mass invariance gives $\int_BA_\chi^m=b^m$. Positivity excludes any residual mass on the extreme set. Thus $\chi/b$ satisfies the same normalized closed-base equation, and \cref{uni-thm:GL-intro} proves the asserted uniqueness.
\end{proof}

The statement selects the limit of the constructed base potentials. It does not identify arbitrary endpoint solutions or give a pointwise limit of the balanced metrics $\gamma_j=L_j^{-1/m}(b_j\kappa+\ddc\phi_j+L_jw)$.

